Señalemos también el siguiente resultado elemental:

\begin{proposition}[6.12.10]
\label{IV.2.s6.12.10-fr}
Sea $X$ un preesquema localmente noetheriano. El conjunto $U_{R_0}(X)$ es
abierto en $X$. Para que $U_{R_1}(X)$ sea abierto en $X$, es necesario y
suficiente que todo punto maximal $x\in X$ tal que $x\in U_{R_1}(X)$ \textup{(lo
que significa que $X$ es \emph{reducido} en $x$)} sea interior a $U_{R_1}(X)$.
\end{proposition}

\begin{proof}
Por definición (5.8.2), $U_{R_0}(X)$ es el conjunto
$X\setmin\bigcup_\alpha X'_\alpha$, donde $X'_\alpha$ recorre el conjunto de las
componentes irreducibles de $X$ tales que $X$ no es reducido en su punto
genérico; como el conjunto de las componentes irreducibles de $X$ es
localmente finito, $U_{R_0}(X)$ es abierto. En cuanto a $U_{R_1}(X)$,
la condición del enunciado es trivialmente necesaria; demostremos que es
suficiente. Sean $X''_\beta$ las componentes irreducibles de $X$ tales
que, en su punto genérico $x''_\beta$, el preesquema $X$ es reducido; por
hipótesis existe un abierto $U\subset U_{R_1}(X)$ que contiene todos los
$x''_\beta$. Designemos entonces por $Z_\lambda$ las componentes
irreducibles del conjunto cerrado $Z=X\setmin U$ cuyo punto genérico
$z_\lambda$ es tal que $\dim(\mathcal O_{X,z_\lambda})\leq1$ y
$\mathcal O_{X,z_\lambda}$ no es regular. Ningún punto perteneciente a
alguno de los $Z_\lambda$ puede pertenecer
\oldpage[IV-2]{168}
a $U_{R_1}(X)$; recíprocamente, si $x\in Z$ no pertenece a ninguno de los
$Z_\lambda$, entonces, para toda generización $x'$ de $x$, o bien $x'$
pertenece a $U$, o bien se tiene $\dim(\mathcal O_{X,x'})\geq2$, o bien se
tiene $\dim(\mathcal O_{X,x'})=1$ y, como $x'$ no pertenece a ninguno de
los $Z_\lambda$, $\overline{\{x'\}}$ es necesariamente una componente
irreducible de $Z$, de codimensión 1 en $X$, distinta de los
$Z_\lambda$; por tanto, $\mathcal O_{X,x'}$ es regular por definición.
Se concluye que $U_{R_1}(X)=X\setmin\bigcup_\lambda Z_\lambda$; como el conjunto
de los $Z_\lambda$ es localmente finito en $Z$,
$\bigcup_\lambda Z_\lambda$ es cerrado, lo que termina de demostrar que
$U_{R_1}(X)$ es abierto en $X$.
\end{proof}
